% Modernised reading — fonds Alexandre Grothendieck, shelfmark 126
% (pages 1-34, the whole folder).
%
% This is an INTERPRETATION, not a transcription. It re-reads the pages in
% today's notation and vocabulary, states the hypotheses the manuscript leaves
% implicit, and drops the critical apparatus so that the mathematics reads
% straight through. Everything the brackets and underlines carried moves into a
% footnote. It is held to being mathematically correct as it stands; where the
% manuscript is loose or mistaken, the true statement is what appears here and
% the footnote says what the page has. In any dispute of substance the
% transcription governs: whoever cites Grothendieck must cite batch-01.fr.tex
% (pp. 1-20) or batch-02.fr.tex (pp. 21-34).
%
% Pass: Opus 5.5 (claude-opus-5-5), 2026-10-03 — first pass, unchecked against the pages by a human.
% The folder was read whole, its two batches together; this is one document
% for the shelfmark. It was made on Opus 5.5 and has not been measured against
% the reference reading of folder 115 (made on Opus 5). The literature named
% in the footnotes is cited from memory and was not checked against the
% sources.
%
% Scope. Pages 15 (a cover sheet bearing only « Descente fidèlement plate »),
% 16, 24, 25 and 30 (blank, or blank letterhead paper) carry no mathematics.
% Much of the prose of pages 2-7 and 21-23 is illegible in the transcription;
% the arguments there are reconstructed from the formulas and flagged as such.
%
% Spine. Not one argument but five pieces on descent, in this order:
%   pp. 1-3    Čech acyclicity of quasi-coherent sheaves for fpqc coverings
%              (Th. 1, Cor. 1-4, remarks)
%   pp. 4-7    application: a nilpotent thickening of an effective flat
%              equivalence relation is effective, with flat quotient
%   p. 8       a three-line plan of the above
%   pp. 9-14   the topology generated by a pretopology T on C and T' on a full
%              subcategory C' (Prop. 1, Cor. 2); the fpqc topology so obtained
%              (Prop. 2) and its sheaf criterion
%   pp. 15-23  under the cover « Descente fidèlement plate »: the radicial /
%              finite étale factorisation of a finite universally open morphism
%              of constant separable degree, by descent of étale coverings
%   pp. 26-29  generic descent: a descent datum defined and satisfying the
%              cocycle condition on dense opens extends to a true one over a
%              dense open of T
%   pp. 31-34  descent along a faithfully flat non-finite morphism: a Galois
%              covering with all but one point of a fibre removed
%
% Notation fixed for the folder. The fpqc topology is called so (he writes
% « topologie fidèlement plate quasi-compacte » or nothing). Čech groups
% \check{H}; \mathcal{H}^q(F) the presheaf T |-> H^q(T, F_T). In pp. 4-7 the
% index 0 marks reductions (X_0 ⊂ X); in pp. 17-23 the étale quotient is
% written X^{ét} (his X_0^*, X_0, X'_0), so as not to collide with that index.
% In pp. 26-29 phi_{ba} : X_a -> X_b is his pointwise notation, kept and
% explained.
%
% Substantive departures, each footnoted where it occurs:
%   - p. 1: F is a sheaf on S (the page has an uncertain « X »); Cor. 1's
%     « X./X » is read X./S;
%   - pp. 4-7: the hypotheses (largely illegible) are reconstructed: p_1, p_2
%     flat, R_0 = R x_X X_0; the sign of lambda is fixed;
%   - p. 9: condition (c), struck on the page but used on p. 13, is restored
%     in a reconstructed form;
%   - p. 17: the remark that « d constant » can be dropped for Y connected is
%     not true as it stands (counterexample in a footnote);
%   - p. 21: Lemma 2's « degré 1 sur Y » is read « sur Y' »;
%   - p. 26: X is in F_T, not F_S; U, V are taken open from the start; the
%     shrinkings are reordered so that V need not be re-shrunk;
%   - p. 28: « W x_{TxT} W ≃ T^4 » is an open subset of T^4, not T^4; the
%     list of conditions defining W''' is corrected to match the computation
%     of p. 29 (t'' and t''' exchanged);
%   - p. 31-32: T x_S T is G x T (|G| copies), not « Delta_T x (G x G) »,
%     and T^3 is G x G x T; the computation of g^{-1}(T') ∩ T' assumes the
%     decomposition group of t trivial, which is added as a hypothesis.
%
% Announced and not established: Remark 1 of p. 3 (illegible); the end of the
% proof of the factorisation (pp. 18-23 are sketches); « À formuler » and the
% « Question » of p. 23; the « Condition de représentation » of p. 29 (stated,
% not proved); the effectivity question of p. 34 (posed, not answered).

\documentclass[11pt,a4paper]{article}
\input{../preamble/grothendieck}

\folder{126}
\pages{1}{34}
\dating{s.d. — le groupe « Descentes » (126 à 133) est daté [avant 1970]}
\watermark{Édition de démonstration\\interprétation personnelle de l'œuvre}
\foldertitle{Descente fidèlement plate. Divers : notes manuscrites (s.d.) — lecture modernisée du dossier entier}

\begin{document}

\section*{Résumé}

\begin{resume}
Ces feuillets cherchent à savoir quand un objet construit « après une
extension » provient en fait d'un objet défini avant elle.

L'image la plus simple est celle des nombres complexes. Un espace vectoriel
complexe muni d'une conjugaison — une application qui se comporte comme
$z \mapsto \bar z$ — est toujours le complexifié d'un espace vectoriel réel :
celui des vecteurs que la conjugaison fixe. La conjugaison est une \emph{donnée
de descente}, et l'espace réel est ce qu'on récupère en « descendant » de
$\mathbf{C}$ à $\mathbf{R}$. La géométrie algébrique a besoin de la même
opération dans une généralité bien plus grande : on remplace un espace $S$ par
un espace $S'$ plus commode qui le recouvre — une réunion de morceaux, une
extension de corps, un revêtement — on y construit ce qu'on veut, et l'on se
demande si la construction redescend. Le mot « fidèlement plat » désigne la
classe très large de recouvrements $S' \to S$ pour lesquels Grothendieck avait
montré que les modules descendent toujours.

Le dossier est fait de cinq pièces sur ce thème, sans autre lien entre elles.
La première calcule, pour un tel recouvrement, l'obstruction à descendre : elle
est mesurée par une cohomologie (celle de Čech), et le théorème dit qu'elle
est nulle pour les modules ; on en tire que les invariants cohomologiques
usuels ne dépendent pas de la finesse des recouvrements autorisés. Une
application suit : si l'on sait former le quotient d'un espace par une
relation d'équivalence, on sait aussi le former après une petite perturbation
infinitésimale de l'espace et de la relation. La deuxième pièce est de pure
logique des recouvrements : elle montre comment fabriquer une notion de
« recouvrement » à partir de deux autres, et retrouve ainsi la topologie
fidèlement plate. La troisième décompose une application finie entre espaces
en une partie qui ne fait que « tasser » les points et une partie qui est un
vrai revêtement, et le prouve en descendant des revêtements. La quatrième
montre qu'une donnée de descente qui n'est connue que « presque partout »
peut être prolongée, quitte à rétrécir l'espace. La dernière examine ce qui se
passe quand le recouvrement n'est plus fini — un revêtement auquel on a ôté
tous les points d'une fibre sauf un — et ramène la question à un problème de
prolongement à travers un point, qu'elle laisse ouvert.

Ce qui frappe, c'est la manière de travailler. Les énoncés sont posés vite,
puis éprouvés sur un cas : la topologie engendrée est essayée aussitôt sur les
schémas, la descente non finie sur un revêtement galoisien découpé à la main.
Les questions sont notées au moment où elles se présentent, avec leurs points
d'interrogation (« Caractérisation des foncteurs représentables par des
revêtements étales ??? »), et une démonstration reçoit un « sauf erreur »
plutôt qu'une certitude affectée. Une feuille de plan de trois lignes
s'intercale au milieu, comme pour fixer où l'on va.

Pour aller plus loin, les noms d'aujourd'hui : descente fidèlement plate et
topologie fpqc, cohomologie de Čech et suite spectrale de Čech, déformations
infinitésimales et quotients plats, topologies de Grothendieck engendrées,
revêtements étales et morphismes radiciels, descente galoisienne, effectivité
de la descente.

\keywords{faithfully flat descent, fpqc topology, Čech cohomology,
Čech-to-derived spectral sequence, quasi-coherent sheaf, nilpotent thickening,
flat equivalence relation, Grothendieck topology, finite étale cover,
radicial morphism, universally open morphism, generic descent, birational
group law, Galois descent, effective descent}
\end{resume}

\pagerange{1}{34}
\section*{Le fil du dossier, et les conventions}

Le dossier compte trente-quatre pages. Les pages 16, 24, 25 et 30 sont
blanches, ou ne portent qu'un en-tête imprimé ; la page 15 est une chemise qui
ne porte que les mots « Descente fidèlement plate ». Les pages 17 à 23 sont
écrites sur du papier à en-tête d'un dactylographe, qui n'a rien à voir avec
leur contenu. Rien ne date les feuillets.

Il n'y a pas de fil unique, mais cinq pièces :

\begin{itemize}
\item \textbf{Cohomologie de Čech fidèlement plate} — pages 1 à 3 : le
Théorème 1 (acyclicité du complexe de Čech d'un faisceau quasi-cohérent pour
un recouvrement fpqc) et ses quatre corollaires.
\item \textbf{Déformation d'une relation d'équivalence plate} — pages 4 à 7,
avec la feuille de plan de la page 8.
\item \textbf{Topologies engendrées} — pages 9 à 14 : une topologie engendrée
par une prétopologie sur une catégorie et une autre sur une sous-catégorie
pleine (Proposition 1, Corollaire 2), et le cas des schémas, qui donne la
topologie fpqc (Proposition 2). La démonstration de la Proposition 1 est
reprise page 13, après l'énoncé de la Proposition 2.
\item \textbf{La factorisation radicielle-étale} — pages 17 à 23, sous la
chemise de la page 15 : un morphisme fini universellement ouvert de degré
séparable constant se factorise en un morphisme radiciel suivi d'un revêtement
étale. Démonstrations esquissées seulement.
\item \textbf{Deux questions de descente} — pages 26 à 29 : la descente
« générique » ; pages 31 à 34 : la descente le long d'un morphisme
fidèlement plat non fini.
\end{itemize}

\emph{Conventions.} « Schéma » a le sens d'aujourd'hui. On appelle
\emph{fpqc} un morphisme fidèlement plat et quasi-compact. Pour un morphisme
$S' \to S$, on note $S'^{(n)} = S' \times_S \cdots \times_S S'$ ($n$
facteurs). Pour un faisceau quasi-cohérent $F$ sur $S$ et un $S$-schéma $T$,
$F_T$ est l'image inverse de $F$ sur $T$ ; $\check{H}^p(S'/S, G)$ désigne la
cohomologie du complexe de Čech
$G(S') \to G(S'^{(2)}) \to G(S'^{(3)}) \to \cdots$ d'un préfaisceau $G$, et
$\mathcal{H}^q(F)$ le préfaisceau $T \mapsto H^q(T, F_T)$. L'indice $0$ ne sert
qu'aux pages 4 à 7, pour les réductions ($X_0 \subset X$) ; le revêtement
étale des pages 17 à 23 est noté $X^{\text{ét}}$ pour qu'on ne le confonde
pas avec elles.\footnote{Grothendieck écrit $X^{*}_0$ à la page 17 (l'exposant
est d'une lecture incertaine), $X_0$ et $X'^{*}_0$ aux pages 21 à 23.}

\pagerange{1}{3}
\section*{Cohomologie de Čech et descente fidèlement plate (pages 1 à 3)}

\textbf{Théorème 1.} Soient $S$ un schéma affine, $f : X \to S$ un morphisme
fidèlement plat, quasi-compact et quasi-séparé, et $F$ un faisceau
quasi-cohérent sur $S$. Alors
\[
  \check{H}^i(X/S, F) \simeq
  \begin{cases}
    H^0(S, F) & \text{si } i = 0, \\
    0 & \text{si } i \neq 0 ,
  \end{cases}
\]
où le complexe de Čech est celui des $\Gamma(X^{(n)}, F_{X^{(n)}})$.%
\footnote{La page lit « $F$ faisceau quasi-cohérent sur $X$ », le $X$ étant
incertain ; le faisceau est nécessairement sur $S$, comme le montre $H^0(S,
F)$. « Quasi-séparé, $S$ affine » est ajouté en interligne. À droite de
l'énoncé, deux lignes encadrées (« couvrant quasi-séparé / quasi-[…] plat »)
semblent proposer une autre rédaction des hypothèses.}

La démonstration n'est pas écrite, mais le Corollaire 1 en donne la marche :
la formation des groupes $\check{H}^i(X/S, F)$ commute à tout changement de
base plat $S_1 \to S$ avec $S_1$ affine (c'est le changement de base plat en
cohomologie des schémas quasi-compacts et quasi-séparés), et un changement de
base fidèlement plat ne change pas l'exactitude d'un complexe. On choisit
$S_1$ égal à la somme d'un nombre fini d'ouverts affines recouvrant $X$ ;
$S_1$ est affine, fidèlement plat sur $S$, et $X \times_S S_1 \to S_1$ a une
section. Or pour un recouvrement qui a une section, le complexe de Čech
augmenté est homotopiquement nul, « et c'est immédiat ».

\textbf{Corollaire 1} (forme faisceautique). Sous les mêmes hypothèses, sans
supposer $S$ affine, les faisceaux $\underline{H}^i(X/S, F)$ sur $S$ — obtenus
en faisant la construction précédente au-dessus des ouverts affines de $S$ —
sont quasi-cohérents, nuls pour $i \neq 0$, et $F \to \underline{H}^0(X/S, F)$
est un isomorphisme.\footnote{La page écrit « faisceaux qu. coh. sur $X$ »
(incertain) et « $\underline{H}^0(X^{\cdot}/X, F)$ » ; on lit $S$ dans les deux
cas.}

\textbf{Corollaire 2.} Pour tout $S$ et tout faisceau quasi-cohérent $F$ sur
$S$, le foncteur $S' \mapsto \Gamma(S', F_{S'})$ sur les $S$-schémas est un
faisceau pour la topologie fpqc. En effet, on est ramené à deux
vérifications, l'une pour les recouvrements de Zariski, l'autre pour un seul
morphisme fpqc entre schémas affines, et celle-ci est le degré $0$ du
Théorème 1.\footnote{Ce corollaire est écrit en travers de la marge gauche,
encadré ; dans l'angle, « compléter […] plus loin ». La réduction à deux types
de recouvrements est justement l'objet de la Proposition 2 des pages 11 à 14.
C'est le théorème de descente fidèlement plate des sections, que Grothendieck
avait publié dans « Technique de descente et théorèmes d'existence en
géométrie algébrique I », Séminaire Bourbaki, exposé 190 (1959) — cité de
mémoire.}

\textbf{Corollaire 3.} Munissons la catégorie des schémas d'une topologie
$\tau$ comprise entre celle de Zariski et la topologie fpqc. Alors pour tout
schéma $X$ et tout faisceau quasi-cohérent $F$ sur $X$, notant $F^{\tau}$ le
faisceau $T \mapsto \Gamma(T, F_T)$ sur le gros site de $X$ (c'en est un par
le Corollaire 2), on a
\[
  H^{*}(X_{\tau}, F^{\tau}) = H^{*}(X, F) .
\]
\emph{Démonstration.} Par la suite spectrale de Leray du morphisme de sites
$X_{\tau} \to X_{\mathrm{Zar}}$, on est ramené à prouver que pour $X$ affine,
$H^i(X_{\tau}, F^{\tau}) = 0$ pour $i > 0$, l'égalité en degré $0$ résultant de
la définition de $F^{\tau}$. Pour cela, le critère de Cartan : il suffit que
la cohomologie de Čech de $F^{\tau}$ soit nulle en degrés $> 0$ pour un système
cofinal de recouvrements des schémas affines, et les recouvrements finis par
des affines fidèlement plats en forment un ; c'est le Théorème 1.%
\footnote{La démonstration, à partir de « Pour la première », est en grande
partie illisible ; on n'en lit que « suite spectrale », « Čech », « pour tout
recouvrement $\mathfrak{X}$ […] fid. plat », « $\check{H}^p(\mathfrak{X}, F) =
0$ pour $p \neq 0$ » et « Th.~1 ». La forme donnée ici est la démonstration
classique, que ces mots suffisent à reconnaître. L'énoncé suppose, comme il
est d'usage, que dans la topologie fpqc les recouvrements d'un schéma affine
sont raffinés par des familles finies de schémas affines.}

\textbf{Corollaire 4.} Soient $S' \to S$ un recouvrement fpqc et $F$ un
faisceau quasi-cohérent sur $S$. Il y a une suite spectrale (de Čech vers la
cohomologie dérivée)
\[
  E_2^{p,q} = \check{H}^p\bigl(S'/S, \mathcal{H}^q(F)\bigr) \Longrightarrow
  H^{p+q}(S, F) ,
\]
où, par le Corollaire 3, peu importe que la cohomologie $H^q(T, F_T)$ soit
prise pour la topologie de Zariski ou pour la topologie fpqc. On en tire la
suite exacte des termes de bas degré :
\[
  0 \to \check{H}^1(S'/S, F) \to H^1(S, F) \to
  \check{H}^0\bigl(S'/S, \mathcal{H}^1(F)\bigr) \to \check{H}^2(S'/S, F) \to
  H^2(S, F) ,
\]
où l'isomorphisme $H^0(S, F) \simeq \check{H}^0(S'/S, F)$ est déjà connu.

\textbf{Remarques.} 1) La première remarque est illisible.\footnote{On n'en lit
que « Bien entendu, on peut […] ».}

2) Si $S$ est affine et $S' \to S$ fpqc et quasi-séparé, la démonstration du
Théorème 1 montre plus : pour tout $q \neq 0$ et tout $p$,
\[
  \check{H}^p\bigl(S'/S, \mathcal{H}^q(F)\bigr) = 0 ;
\]
en particulier
\[
  \check{H}^0\bigl(S'/S, \mathcal{H}^q(F)\bigr) =
  \mathrm{Ker}\bigl(H^q(S', F) \rightrightarrows H^q(S' \times_S S', F)\bigr)
  = 0 .
\]
C'est en effet le même argument : après le changement de base $S_1 \to S$
utilisé plus haut, le complexe $p \mapsto H^q(S'^{(p+1)}, F)$ devient
homotopiquement nul, d'augmentation $H^q(S_1, F_{S_1}) = 0$ puisque $S_1$ est
affine.\footnote{La justification est nôtre ; la page n'en lit que « La
démonstration du Th.~1 montre ». Pour $q = 1$, on peut aussi la lire sur la
suite exacte ci-dessus : $H^1(S, F) = 0$ et $\check{H}^2(S'/S, F) = 0$ quand $S$
est affine. C'est sous cette forme, en degré $1$, que la remarque sert à la
page 6.}

\pagerange{4}{8}
\section*{Déformation d'une relation d'équivalence plate (pages 4 à 8)}

Les hypothèses de cette application sont en grande partie illisibles, et l'on
en donne une forme qui rend correct ce qui suit.\footnote{On lit sur la page :
« On suppose que […] $R_0$ est (M) effective, […] et que $p_1^0 : R_0 \to
X_0$ [est] fid. plat quasi-compact et quasi-séparé. Sous ces conditions, $R$
est M-effective, et si $X_0 \to Y_0$ est plat i.e. $R_0 \to X_0$ […] ». Le
sigle « (M) » n'est pas résolu ; il peut désigner la catégorie dans laquelle
on prend le quotient. La platitude de $p_1, p_2 : R \to X$ et l'égalité $R_0 =
R \times_X X_0$ ne se lisent pas mais sont nécessaires : sans elles l'idéal de
$R_0$ dans $R$ n'est pas l'image inverse de celui de $X_0$ dans $X$, et
l'isomorphisme final de la page 7 tombe.}

\textbf{Proposition.} Soient $X$ un schéma, $X_0 \subset X$ un sous-schéma
fermé défini par un idéal nilpotent $\mathcal{J}_X$, et $p_1, p_2 : R
\rightrightarrows X$ une relation d'équivalence plate (les $p_i$ plats) ;
posons $R_0 = R \times_X X_0$ (par l'une ou l'autre projection), de sorte que
$R_0 \rightrightarrows X_0$ est une relation d'équivalence. Supposons-la
\emph{effective} au sens fpqc : il existe $Y_0$ et $X_0 \to Y_0$ fpqc et
quasi-séparé avec $R_0 \simeq X_0 \times_{Y_0} X_0$. Alors $R$ est effective :
il existe un schéma $Y$, épaississement nilpotent de $Y_0$, et un morphisme
$X \to Y$ fidèlement plat prolongeant $X_0 \to Y_0$, avec $R \simeq X \times_Y
X$.

On peut supposer $\mathcal{J}_X^2 = 0$ : le cas général s'en déduit en
remontant la filtration par les puissances de l'idéal. Alors $\mathcal{J}_X$
est un $\mathcal{O}_{X_0}$-module quasi-cohérent ; notons $\mathcal{J}_R$
l'idéal de $R_0$ dans $R$. Comme les $p_i$ sont plats, $p_i^{*}\mathcal{J}_X
\simeq \mathcal{J}_R$, et ces deux isomorphismes forment une donnée de descente
sur $\mathcal{J}_X$ relativement à $X_0 \to Y_0$. Par descente fidèlement plate
des modules, $\mathcal{J}_X \simeq \mathcal{J}_Y \otimes_{\mathcal{O}_{Y_0}}
\mathcal{O}_{X_0}$ pour un $\mathcal{O}_{Y_0}$-module quasi-cohérent
$\mathcal{J}_Y$.\footnote{La page écrit « $\mathcal{J}_{Y_0}$ faisceau
cohérent sur $Y_0$ » ; « quasi-cohérent » suffit et est ce que donne la
descente sans hypothèse noethérienne.}

\emph{Construction de $Y$.} Sur l'espace topologique de $Y_0$ — qui est aussi,
au-dessus de chaque ouvert, celui de $X$ modulo $R$ — on pose
\[
  \mathcal{A} = \mathrm{Ker}\bigl(\mathcal{O}_X \rightrightarrows
  \mathcal{O}_R\bigr)
\]
(images directes sur $Y_0$ sous-entendues), et l'on veut montrer que la suite
\[
  0 \to \mathcal{J}_Y \to \mathcal{A} \to \mathcal{O}_{Y_0} \to 0
\]
est exacte. Le problème est local sur $Y_0$, qu'on suppose affine. L'exactitude
à gauche et au milieu est la descente des sections : le noyau de $\mathcal{A}
\to \mathcal{O}_{Y_0}$ est $\mathrm{Ker}\bigl(H^0(X_0, \mathcal{J}_X)
\rightrightarrows H^0(R_0, \mathcal{J}_R)\bigr) = H^0(Y_0, \mathcal{J}_Y)$ par
le Théorème 1 en degré $0$, et $\mathrm{Ker}\bigl(H^0(X_0, \mathcal{O}_{X_0})
\rightrightarrows H^0(R_0, \mathcal{O}_{R_0})\bigr) = H^0(Y_0,
\mathcal{O}_{Y_0})$ de même. Reste la surjectivité de $\mathcal{A} \to
\mathcal{O}_{Y_0}$.

\emph{Surjectivité.} Soit $\varphi_0 \in H^0(Y_0, \mathcal{O}_{Y_0})$, et
$f_0 \in H^0(X_0, \mathcal{O}_{X_0})$ son image, qui vérifie $p_1^{*}f_0 =
p_2^{*}f_0$. La suite exacte $0 \to \mathcal{J}_X \to \mathcal{O}_X \to
\mathcal{O}_{X_0} \to 0$ donne une obstruction $\xi \in H^1(X_0,
\mathcal{J}_X)$ à relever $f_0$ en une section de $\mathcal{O}_X$. Ses deux
images par $p_1^{*}$ et $p_2^{*}$ dans $H^1(R_0, \mathcal{J}_R)$ sont les
obstructions à relever $p_1^{*}f_0$ et $p_2^{*}f_0$ en sections de
$\mathcal{O}_R$ ; ces deux fonctions sont égales, donc
\[
  \xi \in \mathrm{Ker}\bigl(H^1(X_0, \mathcal{J}_X) \rightrightarrows
  H^1(R_0, \mathcal{J}_R)\bigr)
  = \check{H}^0\bigl(X_0/Y_0, \mathcal{H}^1(\mathcal{J}_Y)\bigr) ,
\]
qui est nul par la Remarque 2 de la page 3 ($Y_0$ affine, $X_0 \to Y_0$ fpqc et
quasi-séparé). Ainsi $f_0$ se relève en $f \in H^0(X, \mathcal{O}_X)$.

En général $p_1^{*}f \neq p_2^{*}f$, mais leur différence $y = p_2^{*}f -
p_1^{*}f$ est dans $H^0(R_0, \mathcal{J}_R) = H^0(X_0 \times_{Y_0} X_0,
\mathcal{J}_Y)$ : c'est une $1$-cochaîne du complexe de Čech de $X_0/Y_0$ à
coefficients dans $\mathcal{J}_Y$, et c'est un cocycle, $dy = 0$, comme on le
voit en calculant sur $R \times_X R$, dont la réduction est $X_0 \times_{Y_0}
X_0 \times_{Y_0} X_0$. Par le Théorème 1 en degré $1$, $y = d\lambda$ avec
$\lambda \in H^0(X_0, \mathcal{J}_X)$, c'est-à-dire $p_2^{*}\lambda -
p_1^{*}\lambda = y$. Alors $f - \lambda$ relève encore $f_0$ (car $\lambda \in
\mathcal{J}_X$) et vérifie $p_1^{*}(f - \lambda) = p_2^{*}(f - \lambda)$ : c'est
une section de $\mathcal{A}$ qui relève $\varphi_0$.\footnote{La page cherche un
$\lambda$ avec $p_1(f + \lambda) = p_2(f + \lambda)$ et écrit $p_2(\lambda) -
p_1(\lambda) = -y$ (le signe moins est incertain), puis « i.e. $d(\lambda) =
y$ » ; les deux ne sont compatibles qu'avec une convention de signe sur $d$. On
a fixé $d\lambda = p_2^{*}\lambda - p_1^{*}\lambda$ et corrigé le terme
ajouté en conséquence. Le diagramme de la page 6 met en regard les suites
$A \to H^0(X, \mathcal{O}_X) \rightrightarrows H^0(R, \mathcal{O}_R)$ et leurs
réductions.}

\emph{Fin de la démonstration} (page 7). L'espace annelé $Y = (Y_0,
\mathcal{A})$ est un épaississement de $Y_0$ d'idéal quasi-cohérent de carré
nul $\mathcal{J}_Y$ ; c'est donc un schéma.\footnote{Ce point est seulement
indiqué : « $Y_0$ […] défini par un [idéal] nilpotent isomorphe :
$\mathcal{J}_{Y_0}$ ». Qu'un épaississement d'ordre un d'un schéma par un
idéal quasi-cohérent soit un schéma est un énoncé classique, qu'on ne
démontre pas ici.} On a $\mathcal{J}_Y \mathcal{O}_X = \mathcal{J}_X \simeq
\mathcal{J}_Y \otimes_{\mathcal{O}_{Y_0}} \mathcal{O}_{X_0}$ et $X_0$ est plat
sur $Y_0$ ; par le critère de platitude pour les épaississements nilpotents
(le critère local de platitude), $X$ est \emph{plat} sur $Y$. Alors $R$ et $P
= X \times_Y X$ sont plats sur $X$ (par la première projection), le morphisme
$R \to P$ se réduit en l'isomorphisme $R_0 \simeq P_0 = X_0 \times_{Y_0} X_0$,
donc est un isomorphisme.

\begin{tikzcd}
Y & X \arrow[l] & R \arrow[l, bend left=8] \arrow[l, bend right=8] \\
Y_0 \arrow[u, hook] & X_0 \arrow[l] \arrow[u, hook] & R_0 \arrow[l, bend left=8] \arrow[l, bend right=8] \arrow[u, hook]
\end{tikzcd}

Le diagramme est celui de la page 7, où l'on a orienté les flèches
verticales : ce sont les immersions fermées des réductions.\footnote{Sur la
page, le sens de la verticale de gauche, repassée, n'est pas sûr, et celle du
milieu n'a pas de pointe ; la flèche $X_0 \to Y_0$ y est dessinée comme une
inclusion. On les a orientées comme le veut la construction.}

La feuille de la page 8 ne porte que trois lignes de plan : « Généralités sur
la topologie fidèlement plate. Les $\check{H}^p(S'/S, \mathcal{H}^q(F))$ pour
$F$ quasi-cohérent. Applications aux relations d'équivalence… » — c'est le
plan des pages 1 à 7.\footnote{Rien n'assure que la feuille se rapporte aux
pages qui la précèdent plutôt qu'à celles qui suivent ; mais ses deux
dernières lignes décrivent exactement les pages 1 à 7.}

Aujourd'hui, la Proposition se range parmi les énoncés de déformation des
quotients plats : une relation d'équivalence plate qui a un quotient fpqc le
garde sous une déformation infinitésimale plate. La page ne cite rien.

\pagerange{9}{14}
\section*{Topologies engendrées (pages 9 à 14)}

\subsection*{La topologie engendrée par deux prétopologies}

Soient $\mathcal{C}$ une catégorie à produits fibrés munie d'une
prétopologie $T$, et $\mathcal{C}'$ une sous-catégorie pleine munie d'une
prétopologie $T'$, avec les conditions :

(a) tout objet $X$ de $\mathcal{C}$ a une famille $T$-couvrante $(U_\alpha \to
X)$ formée d'objets de $\mathcal{C}'$ ;

(b) $\mathcal{C}'$ est stable par isomorphismes et par produits fibrés dans
$\mathcal{C}$ ;

(c) si $U \in \mathcal{C}'$, $(U_i \to U)$ est $T'$-couvrante et, pour chaque
$i$, $(U_{ij} \to U_i)$ est $T$-couvrante avec les $U_{ij}$ dans
$\mathcal{C}'$, alors la famille composée $(U_{ij} \to U)$ a une sous-famille
$T'$-couvrante.\footnote{La condition (c) est entièrement biffée à la page 9,
de grandes croix et de traits ondulés, et sa lecture est très incertaine ;
mais elle est invoquée à la page 13 (« Par (c), […] extraire de
$U_{\alpha\beta\gamma} \to U_\alpha$ une famille couvrante par $T'$ »), et une
note de marge de la page 10 semble dire que la proposition en a besoin. La
forme donnée ici est reconstituée à partir de cet usage et de la
« situation » dessinée page 9 ($U \leftarrow U_i \leftarrow U_{ij} \leftarrow
U_{ij\mu}$, $T'$ puis $T$ puis $T'$, d'où l'on veut extraire une famille
$T'$-couvrante). Prise avec la famille $T'$-couvrante triviale $(U \to U)$,
elle dit en particulier qu'une famille $T$-couvrante d'un objet de
$\mathcal{C}'$ par des objets de $\mathcal{C}'$ contient une famille
$T'$-couvrante.}

Soit $\mathcal{T}$ la topologie sur $\mathcal{C}$ engendrée par $T$ et $T'$.

\textbf{Proposition 1.} Une famille $(X_i \to X)$ de $\mathcal{C}$ est
$\mathcal{T}$-couvrante si et seulement s'il existe une famille $T$-couvrante
$(U_\alpha \to X)$ avec les $U_\alpha$ dans $\mathcal{C}'$, pour chaque
$\alpha$ une famille $T'$-couvrante $(U_{\alpha\beta} \to U_\alpha)$, et pour
chaque $(\alpha, \beta)$ un indice $i = i(\alpha, \beta)$ et un
$X$-morphisme $U_{\alpha\beta} \to X_i$ :

\begin{tikzcd}[column sep=small, row sep=small]
X_i \arrow[d] & U_{\alpha\beta} \arrow[l] \arrow[d] \\
X & U_\alpha \arrow[l]
\end{tikzcd}

Si $X$ est dans $\mathcal{C}'$, il faut et il suffit aussi qu'il
existe une famille $T'$-couvrante $(U_\alpha \to X)$ et pour chaque $\alpha$
un $X$-morphisme $U_\alpha \to X_{i(\alpha)}$.

\emph{Démonstration.} Si la condition est remplie, $(U_{\alpha\beta} \to X)$
est $\mathcal{T}$-couvrante par transitivité, donc $(X_i \to X)$ aussi par
saturation. Réciproquement, il suffit de voir que les familles du type décrit
forment une topologie : stabilité par isomorphisme, par changement de base et
par composition.\footnote{La page parle des « conditions 1), 2), 3), 4) » sans
les énoncer ; elle dit 4) triviale et 3) « aussi », et traite 1), la
transitivité, en dernier. La démonstration, commencée page 10, est achevée
page 11 (« ok ») puis reprise plus au long page 13, qu'on suit ici.}

Pour un isomorphisme $X' \to X$, on prend par (a) un recouvrement $(U_\alpha
\to X)$ par des objets de $\mathcal{C}'$, les $U_{\alpha\beta}$ réduits à
$U_\alpha$ et les flèches $U_\alpha \to X'$ déduites.

Changement de base par $X' \to X$ (page 13) : les $U'_\alpha = U_\alpha
\times_X X'$ recouvrent $X'$ pour $T$ ; on recouvre chacun pour $T$ par des
$U'_{\alpha\lambda}$ dans $\mathcal{C}'$, ce que permet (a) ; les
$U'_{\alpha\lambda}$ recouvrent $X'$ pour $T$ par transitivité, et les
$U'_{\alpha\beta\lambda} = U_{\alpha\beta} \times_{U_\alpha}
U'_{\alpha\lambda}$, qui sont dans $\mathcal{C}'$ par (b), recouvrent
$U'_{\alpha\lambda}$ pour $T'$ par changement de base.

Transitivité (pages 13 et 11) : soient $(X_i \to X)$ comme ci-dessus, avec
$U_\alpha$, $U_{\alpha\beta}$, $i(\alpha, \beta)$, et pour chaque $i$ une
famille $(X_{ij} \to X_i)$ du même type. Par changement de base, $(X_{ij}
\times_{X_i} U_{\alpha\beta} \to U_{\alpha\beta})_j$ est du même type : il y
a des $U_{\alpha\beta\gamma} \to U_{\alpha\beta}$ $T$-couvrants dans
$\mathcal{C}'$, des $U_{\alpha\beta\gamma\delta} \to U_{\alpha\beta\gamma}$
$T'$-couvrants, et des morphismes $U_{\alpha\beta\gamma\delta} \to X_{ij}$.
Par (c), on extrait de $(U_{\alpha\beta\gamma} \to U_\alpha)_{\beta, \gamma}$
une famille $T'$-couvrante $(V_{\alpha\lambda} \to U_\alpha)$ ; les
$U_{\alpha\beta\gamma\delta}$ correspondants forment, par composition dans
$T'$, une famille $T'$-couvrante de $U_\alpha$, dont chaque membre s'envoie
dans un $X_{ij}$. C'est ce qu'il fallait.

Le second énoncé, pour $X$ dans $\mathcal{C}'$, résulte du premier et de (c).

\textbf{Corollaire 2.} Un préfaisceau $F : \mathcal{C}^{\circ} \to
(\mathrm{Ens})$ est un faisceau pour $\mathcal{T}$ si et seulement si c'est un
faisceau pour $T$ et que sa restriction à $\mathcal{C}'$ est un faisceau pour
$T'$.\footnote{La justification, en marge, est illisible sauf « OK ». Le point
à voir est que la condition de faisceau pour les changements de base hors de
$\mathcal{C}'$ des familles $T'$-couvrantes se ramène, par (a), (b) et la
condition de faisceau pour $T$, à celle sur $\mathcal{C}'$.}

\pagerange{11}{14}
\subsection*{Le cas des schémas : la topologie fpqc}

\textbf{Proposition 2.} Prenons pour $\mathcal{C}$ la catégorie des schémas,
avec la prétopologie $T$ de Zariski, pour $\mathcal{C}'$ celle des schémas
affines, et pour $T'$ la prétopologie dont les familles couvrantes $(U_i \to
U)$ sont les familles \emph{finies} de morphismes \emph{plats} dont les images
recouvrent $U$. Les conditions (a), (b), (c) sont remplies.\footnote{Pour (c) :
un recouvrement de Zariski d'un affine par des affines a une sous-famille
finie, et une composée finie de morphismes plats conjointement surjectifs en
est encore une. La page ne vérifie pas les conditions. En marge, en travers,
une liste de variantes de $T'$ : « plate de présentation finie, plate […],
étale » — qui donneraient de même les topologies fppf et étale.} Soit
$\mathcal{T}$ la topologie engendrée. Alors :

(i) Une famille $(X_i \to X)$ est $\mathcal{T}$-couvrante si et seulement si
$X$ est réunion d'ouverts affines $U_\alpha$ et que pour chaque $\alpha$ il
existe une famille finie de morphismes plats $U_{\alpha\beta} \to U_\alpha$,
les $U_{\alpha\beta}$ affines, dont les images recouvrent $U_\alpha$, et qui
se factorisent chacun par un $X_i$.

(ii) Un foncteur $F : (\mathrm{Sch})^{\circ} \to (\mathrm{Ens})$ est un
faisceau pour $\mathcal{T}$ si et seulement si (a) $F$ est un faisceau pour la
topologie de Zariski — « $F$ est de nature locale » — et (b) pour tout
morphisme fidèlement plat $S' \to S$ entre schémas affines, la suite
\[
  F(S) \to F(S') \rightrightarrows F(S' \times_S S')
\]
est exacte.

(iii) Un morphisme fidèlement plat et quasi-compact $X' \to X$ est
$\mathcal{T}$-couvrant ; de même un morphisme fidèlement plat localement de
présentation finie ; plus généralement, une famille de morphismes plats et
\emph{ouverts} dont les images recouvrent $X$ est
$\mathcal{T}$-couvrante.\footnote{Sur la page, (iii) est écrit avant (ii),
auquel il est rattaché par un trait de marge ; on a rétabli l'ordre des
numéros. Pour la dernière assertion : un ouvert affine de $X$ est
quasi-compact, donc recouvert par les images, ouvertes, d'un nombre fini
d'ouverts affines des $X_i$.}

\emph{Démonstration.} (i) est la Proposition 1. (ii) résulte du Corollaire 2 :
la condition (a) fait que $F$ transforme les sommes finies en produits, ce
qui ramène la condition de faisceau pour une famille finie $(U_i \to U)$ de
$T'$ à la condition (b) pour le seul morphisme $\coprod U_i \to U$.%
\footnote{Cette fin de démonstration est à la page 14, après deux lignes
biffées qui continuent une phrase interrompue au bas de la page 12. La
démonstration de (iii) est annoncée (« (iii) par des produits… ») et ne
suit pas.}

La topologie $\mathcal{T}$ est donc la topologie \emph{fpqc}, telle qu'on la
définit aujourd'hui précisément par (i) ; (ii) en est le critère de faisceau
usuel, et le Corollaire 2 de la page 1 dit que les faisceaux quasi-cohérents
en sont.

\pagerange{15}{23}
\section*{Descente fidèlement plate : la factorisation radicielle-étale (pages 15 à 23)}

Ces pages suivent la chemise « Descente fidèlement plate » de la page 15. Elles
s'ouvrent sur un plan à deux entrées — « A) Application à la comparaison entre
les variétés étales [et les variétés] holomorphes » et « B) Application à la
factorisation canonique » — dont seule la seconde est traitée.\footnote{La
lecture « étales » est incertaine, de même que « factorisation ». L'entrée
A) n'a pas de suite dans le dossier.}

\pagerange{17}{20}
\subsection*{L'énoncé, et le revêtement des numérotations}

Pour un morphisme fini $f : X \to Y$ et $y \in Y$, notons $d(y)$ le
\emph{degré séparable} de la fibre $f^{-1}(y)$ sur $k(y)$, c'est-à-dire le
nombre de ses points géométriques.

\textbf{Proposition.} Soit $f : X \to Y$ un morphisme fini et universellement
ouvert, avec $Y$ localement noethérien, et supposons $d(y)$ constant. Alors
$f$ se factorise en
\[
  X \xrightarrow{\ f_1\ } X^{\text{ét}} \xrightarrow{\ f_2\ } Y ,
\]
où $f_1$ est fini, surjectif et radiciel, et $f_2$ un revêtement étale (fini
étale). De plus tout $Y$-morphisme $X \to Z$, avec $Z$ étale sur $Y$, se
factorise de façon unique par $X^{\text{ét}}$.\footnote{« Univ. ouvert » est
ajouté en interligne, d'une lecture incertaine ; c'est lui qui fait marcher
l'argument de la page 19. La page ajoute qu'« on peut supprimer » l'hypothèse
« $d(y)$ constant » quand $Y$ est connexe, « quitte à remplacer $Y$ par une
partie […] fermée » ; la fin est illisible. Telle quelle, l'assertion ne
tient pas : pour $k$ de caractéristique $\neq 2$, $X = \mathrm{Spec}\,
k[t, x]/(x^2 - t) \to Y = \mathrm{Spec}\, k[t]$ est fini et plat, donc
universellement ouvert, $Y$ est connexe, et $d$ vaut $1$ en $t = 0$ et $2$
ailleurs ; aucune factorisation radicielle-étale n'existe. L'hypothèse de
constance est donc à garder. Le dossier n'achève pas la démonstration, et la
présente lecture ne certifie pas l'énoncé dans toute sa généralité.}

Les pages qui suivent en esquissent la démonstration par deux voies.

\emph{L'espace des numérotations} (pages 18 et 20). Soient $n$ la valeur
constante de $d$, $I = \{1, \ldots, n\}$, et
\[
  Y' = X^{I} = X \times_Y \cdots \times_Y X
\]
($n$ facteurs). Les $n$ projections $Y' \to X$ sont $n$ sections du
morphisme $X' = X \times_Y Y' \to Y'$, d'où un $Y'$-morphisme $Y' \times I \to
X'$, et pour chaque $y' \in Y'$ une application $I \times \mathrm{Spec}\,
k(y') \to X'_{y'}$. Soit $Y'' \subset Y'$ l'ensemble des points où elle est
injective sur les points géométriques — donc bijective, puisque la fibre a
$n$ points géométriques : $Y''$ paramètre les \emph{numérotations} des points
géométriques des fibres de $f$.

\textbf{Lemme.} $Y''$ est ouvert dans $Y'$, et $Y'' \to Y$ est surjectif. Un
point $y'$ est dans $Y''$ si et seulement s'il a un voisinage ouvert $U$ tel
que $U \times I \to X'|_U$ soit une immersion fermée.\footnote{La page 18 est
biffée de deux longues diagonales, et sa suite est à la page 20. L'énoncé du
lemme y dit aussi quelque chose de « fermé », qu'on ne sait pas lire, et
« $Y'' \to$ connexe » (incertain). Les justifications sont nôtres : sur un
corps, $n$ points rationnels distincts d'un schéma fini forment un
sous-schéma fermé, de sorte qu'« injectif » équivaut à « immersion fermée »
sur la fibre — la page hésite entre « radiciel », « monomorphisme » et
« immersion fermée », et c'est la dernière qui convient ; par le lemme de
Nakayama, un morphisme fini qui est une immersion fermée sur une fibre en est
une au-dessus d'un voisinage. La surjectivité de $Y'' \to Y$ vient de ce que
toute numérotation d'une fibre géométrique est réalisée par un point de
$Y'$.}

\emph{La section ouverte et fermée} (page 19). Si $f : X \to Y$, fini et
universellement ouvert de degré séparable constant $n$, a une section $s$,
alors $s(Y)$ est ouvert et fermé dans $X$. On se ramène à $Y$ irréductible en
regardant les $f^{-1}(Y_i)$ au-dessus des composantes irréductibles $Y_i$ de
$Y$. Alors $s(Y)$ est une réunion de composantes irréductibles de $X$, puisque
$f$ est fini ; il faut voir qu'elle ne rencontre pas la réunion $X''$ des
autres composantes, ce que la page tire de la constance de $d$ : un point
commun ferait chuter le nombre de points géométriques de la
fibre.\footnote{Ce passage est encadré et barré de deux traits obliques, puis
remplacé par « … et alors les composantes irréductibles de $X$ sont
disjointes [vérif.] ». La justification par la chute de $d$ est la nôtre, et
elle utilise que $X''$ a, au-dessus de chaque point, au moins $n - 1$ points
géométriques, ce qu'assure l'ouverture universelle ; on ne l'a pas écrite en
détail.} D'où une décomposition
\[
  X = X_{(1)} \amalg X^{(1)}, \qquad d(X_{(1)}/Y) = 1, \quad d(X^{(1)}/Y) = n - 1 .
\]

\pagerange{21}{23}
\subsection*{Descente des revêtements étales}

On itère (page 21) : avec $Y_1 = X$, le morphisme $X_1 = X \times_Y Y_1 \to
Y_1$ a la section diagonale, donc $X_1 = X_{(1)} \amalg X^{(1)}$ ; on pose
$Y_2 = X^{(1)}$, et ainsi de suite. On obtient :

\textbf{Lemme 2.} Il existe $Y' \to Y$ fini, surjectif et universellement
ouvert tel que $X' = X \times_Y Y'$ soit somme de parties ouvertes et fermées
$X'_i$ de degré séparable $1$ sur $Y'$.\footnote{La page écrit « de degré […]
séparable 1 sur $Y$ » ; c'est sur $Y'$ que le degré est $1$. L'indice de
$Y'_2$ est écrit dans l'énoncé puis omis. « Radiciel » est biffé devant
« universellement ouvert ». La lecture de toute la prose des pages 21 à 23
est très incertaine : encre pâle, main très rapide.}

Soit $I$ l'ensemble des composantes $X'_i$ ; $Y' \times I$ est un revêtement
étale trivial de $Y'$, et $X' \to Y' \times I$ est radiciel et surjectif. Le
revêtement $Y' \times I$ hérite de $X'$ une donnée de descente relative à $Y'
\to Y$ ; comme les morphismes finis surjectifs sont de descente effective pour
les revêtements étales, elle définit un revêtement étale $X^{\text{ét}}$ de
$Y$ dont $Y' \times I$ est l'image inverse.\footnote{Le fait invoqué est
dans SGA~1, exposé IX, § 4 (cité de mémoire), sous une hypothèse de
présentation finie qu'assure ici le caractère localement noethérien de $Y$. La
page dit seulement « d'où un rev. étale $X$ de $Y$ ».} On choisit ensuite un
revêtement étale surjectif $Y^{*} \to Y$ qui déploie complètement
$X^{\text{ét}}$, et l'on trouve :

\textbf{Lemme 3.} Il existe un morphisme étale surjectif $Y^{*} \to Y$ tel que
les composantes connexes de $X^{*} = X \times_Y Y^{*}$ soient de degré
séparable $1$ sur $Y^{*}$.

Le morphisme canonique $X^{*} \to Y^{*} \times I$ (où $I$ est maintenant
l'ensemble des composantes connexes de $X^{*}$, $Y^{*}$ supposé connexe) est
compatible aux données de descente relatives à $Y^{*} \to Y$, et définit donc
un morphisme $X \to X^{\text{ét}}$, fini, surjectif et radiciel. La
propriété universelle se vérifie de même après le changement de base étale
$Y^{*} \to Y$, où $X^{\text{ét}}$ devient trivial et où elle est
« immédiate ».\footnote{La phrase est presque entièrement illisible ; on n'en
lit que les mots clés (« universel », « $X^{*} \to X^{*}_0 \to Y'^{*}$ »,
« immédiat »). Ce paragraphe dit ce qu'elle doit établir, non ce qu'elle
établit.}

La page 23 se clôt sur deux notes. « À formuler » : un énoncé pour les
morphismes universellement ouverts non nécessairement finis, où il semble que
le degré séparable soit remplacé par la connexité géométrique des
fibres.\footnote{Lecture conjecturale, à partir des seuls mots lisibles :
« univ. ouvert », « connexité géométrique des fibres ».} Et une
\emph{question} : « Caractérisation des foncteurs représentables par des
revêtements étales ??? ».

Les noms d'aujourd'hui sont ceux du revêtement étale associé à un morphisme
fini, ou de sa « partie étale » ; c'est une forme étale de la factorisation de
Stein, et la question finale est celle des faisceaux étales représentables par
des revêtements, que la théorie des espaces algébriques et des faisceaux
localement constants constructibles a depuis posée en général.

\pagerange{26}{29}
\section*{Descente générique (pages 26 à 29)}

\emph{Données.} $S$ est un schéma noethérien, $T \to S$ un morphisme plat de
type fini ; on note $p_i$, $p_{ij}$ les projections de $T \times_S T$ et de
$T \times_S T \times_S T$. $F$ est une catégorie fibrée sur la catégorie des
$S$-schémas de type fini, telle que les morphismes plats surjectifs (de type
fini) soient des morphismes de \emph{$F$-descente} : l'image inverse par un tel
morphisme est pleinement fidèle sur les objets munis de données de
descente.\footnote{La ligne qui énonce l'hypothèse sur $F$ est très
surchargée ; on y lit « on suppose les morphismes […] » et, biffés, « fid.
plats » et « $F$-descente ». C'est à la page 28 que l'hypothèse sert : « $W
\to T \times_S T$ est surjectif, plat, il s'ensuit que c'est un morphisme de
$F$-descente ». « Morphisme de $F$-descente » est pris au sens de SGA~1,
exposé VI : pleine fidélité, sans effectivité ; l'effectivité est
« $F$-descente stricte », qu'utilise la condition de la page 29.} Soient
enfin :

\begin{itemize}
\item $U$ un ouvert dense de $T \times_S T$ et $V$ un ouvert dense de $T
\times_S T \times_S T$, avec $V \subset \bigcap_{i \neq j}
p_{ij}^{-1}(U)$ ;
\item $X$ un objet de $F_T$ et un isomorphisme $\varphi : p_1^{*}X|_U
\xrightarrow{\sim} p_2^{*}X|_U$ qui satisfait à la condition de cocycle
$p_{13}^{*}\varphi = p_{23}^{*}\varphi \circ p_{12}^{*}\varphi$ au-dessus de
$V$.\footnote{La page écrit « $X \in F_S$ », et $\varphi : (X \times_S T)|U
\to (T \times_S X)|U$. Pour un $X$ sur $S$ ces deux objets seraient le même,
et la question vide ; il faut $X$ sur $T$, comme le confirment « $X_i =
p_i^{*}(X)$ » à la page 27 et « $X \simeq Y \times_{S'} T'$ » à la page 29. La
page prend $U$, $V$ constructibles denses, puis remarque qu'on peut les
supposer ouverts : un constructible dense contient un ouvert dense. On les
prend ouverts d'emblée.}
\end{itemize}

\emph{Notation ponctuelle.} Comme la page, on écrit $\varphi_{ba} : X_a \to
X_b$ pour la valeur de $\varphi$ en un point $(a, b)$ de $U$ ; la condition
de cocycle s'écrit $\varphi_{ca} = \varphi_{cb}\varphi_{ba}$ pour $(a, b, c)
\in V$. Chaque égalité ponctuelle ci-dessous abrège une égalité entre images
inverses de $\varphi$ par des projections, au-dessus d'un ouvert d'une
puissance de $T$ ; c'est sous cette forme qu'elle est vraie.

\textbf{Proposition.} Sous ces conditions, il existe un ouvert dense $T'$ de
$T$ et une \emph{donnée de descente} $\varphi'$ sur $X|_{T'}$ relativement à
$T' \to S$, c'est-à-dire un isomorphisme $\varphi' : p_1^{*}X
\xrightarrow{\sim} p_2^{*}X$ sur tout $T' \times_S T'$ satisfaisant à la
condition de cocycle, telle que $\varphi' = \varphi$ au-dessus de $U \cap (T'
\times_S T')$.\footnote{La page dit « partie constructible dense $T'$ de
$T$ » ; l'argument ci-dessous donne un ouvert dense, ce qui est plus fort et
nécessaire pour que $X|_{T'}$ ait un sens.}

\emph{Réductions} (page 26). Par les théorèmes de constructibilité (EGA~IV,
§ 9, cité de mémoire), l'ensemble des $\tau \in U$ tels que $V \cap
p_{ij}^{-1}(\tau)$ soit dense dans la fibre $p_{ij}^{-1}(\tau)$ pour tout $i
\neq j$ est constructible ; il contient les points maximaux de $U$, car $p_{ij}$
est plat. Il contient donc un ouvert dense $U_1 \subset U$. De même, l'ensemble
des $t \in T$ tels que $U_1 \cap p_i^{-1}(t)$ soit dense dans $p_i^{-1}(t)$
pour $i = 1, 2$ contient un ouvert dense $T'$. On remplace $T$ par $T'$, $U$ par
$U_1 \cap (T' \times_S T')$ et $V$ par $V \cap T'^{3}$ pour les constructions
qui suivent, tout en gardant $\varphi$ sur l'ouvert $U$ d'origine, qui contient
toutes les paires extraites de $V$.\footnote{La page remplace $U$ par $U_1$
sans plus ; il faudrait alors rétrécir $V$ pour garder $V \subset \bigcap
p_{ij}^{-1}(U_1)$, ce qui pourrait détruire la densité qu'on vient
d'obtenir. Garder $\varphi$ sur le $U$ d'origine évite ce cercle. La
justification de la page renvoie à une référence illisible (« cf. […] »). Une
ligne biffée proposait de supposer en outre $U$ symétrique.}

Le principe de tous les arguments qui suivent est le même : chacune des
conditions « $(a, b) \in U$ » ou « $(a, b, c) \in V$ » portant sur un point
variable définit, dans la fibre en ce point — qui est $T$ étendu à un corps —,
un ouvert dense, grâce aux réductions et à la platitude ; un nombre fini
d'ouverts denses d'une même fibre non vide se rencontrent ; les morphismes
de projection qu'on construit sont donc surjectifs, et ils sont plats comme
restrictions à des ouverts de projections $T^{m} \to T^{k}$.

\emph{Construction de $\psi$} (page 27). Posons
\[
  W = p_{12}^{-1}(U) \cap p_{23}^{-1}(U) \subset T \times_S T \times_S T ,
  \qquad \pi = p_{13}|_W : W \to T \times_S T .
\]
Un point de $W$ est un triplet $(t_1, t', t_2)$ avec $(t_1, t')$ et $(t', t_2)$
dans $U$ ; $\pi$ est plat et, par les réductions, surjectif. Soient $X_i =
p_i^{*}X$ sur $T \times_S T$. On définit $\psi : \pi^{*}X_1 \xrightarrow{\sim}
\pi^{*}X_2$ comme le composé passant par le facteur du milieu :
\[
  \psi_{(t_1, t', t_2)} = \varphi_{t_2 t'}\,\varphi_{t' t_1} ,
\]
c'est-à-dire $p_1^{(3)*}X \xrightarrow{\sim} p_2^{(3)*}X \xrightarrow{\sim}
p_3^{(3)*}X$.

\emph{$\psi$ descend.} Il faut voir que $\psi$ est compatible à la donnée de
descente canonique relative à $\pi$, c'est-à-dire que ses deux images inverses
sur
\[
  W \times_{T \times_S T} W \subset T^{4} , \qquad \text{points } (t_1, t',
  t'', t_2) ,
\]
coïncident : $\varphi_{t_2 t'}\varphi_{t' t_1} = \varphi_{t_2
t''}\varphi_{t'' t_1}$.\footnote{À la page 28, Grothendieck écrit $W
\times_{T \times_S T} W \simeq T^{4}$, ajouté au-dessus de la ligne ; c'est
seulement un ouvert de $T^{4}$ (les quatre paires $(t_1, t')$, $(t', t_2)$,
$(t_1, t'')$, $(t'', t_2)$ doivent être dans $U$), comme il l'avait écrit
page 27 (« $\subset$ »).} On introduit pour cela le schéma $W' \subset T^{5}$
des $(t_1, t', t'', t_2, t''')$ tels que
\[
  (t_1, t''', t'),\ (t''', t', t_2),\ (t_1, t''', t''),\ (t''', t'', t_2),\
  (t_1, t''', t_2) \in V .
\]
Sur $W'$, la condition de cocycle donne
\[
  \varphi_{t_2 t'}\varphi_{t' t_1}
  = \varphi_{t_2 t'}(\varphi_{t' t'''}\varphi_{t''' t_1})
  = (\varphi_{t_2 t'}\varphi_{t' t'''})\varphi_{t''' t_1}
  = \varphi_{t_2 t'''}\varphi_{t''' t_1} ,
\]
et de même $\varphi_{t_2 t''}\varphi_{t'' t_1} = \varphi_{t_2
t'''}\varphi_{t''' t_1}$ ; d'où l'égalité cherchée au-dessus de
$W'$.\footnote{La page énumère les triplets avec $t'''$ en tête : « $(t''',
t', t_1)$, $(t''', t', t_2)$, $(t''', t'', t_1)$, $(t''', t'', t_2)$,
$(t''', t_1, t_2) \in V$ » ; on les a ordonnés comme l'exige le calcul qu'elle
fait ensuite, la condition de cocycle n'étant supposée que dans un ordre. Le
dernier triplet sert à la comparaison finale de $\varphi$ et $\varphi'$. Le
schéma de la page 27, en étoile autour de $t'''$, figure ces conditions.}
Or $W' \to W \times_{T \times_S T} W$ est plat et surjectif — « sauf
erreur », écrit la page —, donc un morphisme de $F$-descente, et l'égalité
descend. Comme $\pi$ est lui aussi plat et surjectif, $\psi$ provient d'un
unique isomorphisme
\[
  \varphi' : X_1 \xrightarrow{\sim} X_2 \quad \text{sur } T \times_S T ,
  \qquad \varphi'_{t_2 t_1} = \varphi_{t_2 t'}\varphi_{t' t_1} \ \text{pour tout
  } t' \text{ convenable.}
\]

\emph{Condition de cocycle pour $\varphi'$} (pages 28 et 29). Soit $W''' \subset
T^{6}$ le schéma des $(t_1, t_2, t_3, t', t'', t''')$ tels que
\[
  (t_1, t'),\ (t', t_2),\ (t_2, t''),\ (t'', t_3) \in U
  \quad \text{et} \quad
  (t_1, t', t'''),\ (t''', t'', t_3) \in V ;
\]
la projection $W''' \to T^{3}$ est plate et surjective. Sur $W'''$,
\begin{align*}
  \varphi'_{t_3 t_2}\varphi'_{t_2 t_1}
  &= (\varphi_{t_3 t''}\varphi_{t'' t_2})(\varphi_{t_2 t'}\varphi_{t' t_1})
   = \varphi_{t_3 t''}(\varphi_{t'' t'''}\varphi_{t''' t'})\varphi_{t' t_1} \\
  &= (\varphi_{t_3 t''}\varphi_{t'' t'''})(\varphi_{t''' t'}\varphi_{t' t_1})
   = \varphi_{t_3 t'''}\varphi_{t''' t_1} = \varphi'_{t_3 t_1} ,
\end{align*}
où la deuxième égalité vient de ce que $\varphi_{t'' t_2}\varphi_{t_2 t'}$ et
$\varphi_{t'' t'''}\varphi_{t''' t'}$ valent tous deux $\varphi'_{t'' t'}$, et
la quatrième de la condition de cocycle sur les deux triplets de $V$. La
condition de cocycle descend de $W'''$ à $T^{3}$.\footnote{La liste de la
page 28 porte $(t_2, t''')$, $(t''', t_3)$ dans $U$ et $(t'', t''', t_3)$ dans
$V$ : $t''$ et $t'''$ y sont échangés par rapport au calcul de la page 29, que
l'on suit. Les paires $(t', t''')$, $(t''', t'')$, $(t_1, t''')$, $(t''',
t_3)$, dont le calcul a aussi besoin, sont dans $U$ parce qu'elles sont
extraites de triplets de $V$. La phrase qui conclut la page 28 est
illisible.}

\emph{$\varphi'$ prolonge $\varphi$} (page 29). Pour $(t_1, t_2, t_3) \in V$, on
peut prendre $t' = t_2$ dans la définition de $\varphi'$ : $\varphi'_{t_3
t_1} = \varphi_{t_3 t_2}\varphi_{t_2 t_1} = \varphi_{t_3 t_1}$. Ainsi
$p_{13}^{*}\varphi' = p_{13}^{*}\varphi$ sur $V$, et comme $p_{13} : V \to U$
est plat et surjectif, $\varphi' = \varphi$ sur $U$. « La démonstration […]
facile », conclut la page.

\textbf{Condition de représentation.} Supposons que pour tout ouvert dense $T'$
de $T$, il existe un ouvert dense $S'$ de $S$ au-dessus duquel $T'_{S'} \to S'$
soit un morphisme de $F$-descente \emph{stricte} (la descente y est
effective). Alors, quitte à rétrécir $T'$, on trouve $Y \in F_{S'}$ et un
isomorphisme $X|_{T'} \simeq Y \times_{S'} T'$ compatible aux données de
descente.\footnote{L'énoncé n'est pas démontré, et ses quantificateurs sont
en partie illisibles ; la page dit « $S'$ (…) de $S$ ». On a pris pour $S'$
un ouvert dense, ce qui est le sens « générique » de toute la section. Le
« $S$ » indicé de « $F_{S'}$ » est surchargé.}

Le procédé — définir l'objet cherché en un couple de points génériques
$(t_1, t_2)$ en passant par un troisième point $t'$ « en position générale »,
puis vérifier l'indépendance par un quatrième — est celui du théorème de Weil
sur les lois de groupe birationnelles et de son critère de descente pour les
variétés (A. Weil, « The field of definition of a variety », 1956, cité de
mémoire) ; on en trouve une version schématique, due à M. Artin, dans SGA~3,
exposé XVIII. La page ne nomme ni l'un ni l'autre.

\pagerange{31}{34}
\section*{Descente pour un morphisme fidèlement plat non fini (pages 31 à 34)}

Question : la descente est-elle possible le long d'un morphisme fidèlement
plat \emph{non fini} ? Un cas typique est le suivant.

Soient $f : T \to S$ un revêtement étale galoisien de groupe $G$ — un
$G$-torseur fini étale —, $s$ un point fermé de $S$ et $t \in f^{-1}(s)$. On
suppose que le groupe de décomposition de $t$ est trivial, c'est-à-dire que
$G$ permute librement les points de $f^{-1}(s)$ (par exemple si $k(s)$ est
séparablement clos).\footnote{L'hypothèse n'est pas sur la page, mais le
calcul qui suit en dépend : si $g \neq e$ fixe $t$, alors $g^{-1}(T') \cap T'
= T'$ et non $T''$, et la conclusion change.} Soit $T' = T - \Phi$, où $\Phi
= f^{-1}(s) - \{t\}$ : c'est un ouvert, et $T' \to S$ est étale et surjectif,
fidèlement plat, mais il n'est pas fini. Notons $T'' = T - f^{-1}(s)$ et
$S'' = S - \{s\}$.

\emph{Le produit $T' \times_S T'$.} Comme $T$ est un $G$-torseur,
\[
  T \times_S T = \coprod_{g \in G} U_g , \qquad U_g = (g \times e)(\Delta_T)
  \simeq T ,
\]
l'isomorphisme $\varphi_g : T \xrightarrow{\sim} U_g$ étant $\tau \mapsto (g\tau,
\tau)$.\footnote{La page écrit « $T \times_S T = \Delta_T \times (G \times G)$ »,
après avoir biffé un signe de produit indicé $G$. Il y a $|G|$ composantes et
non $|G|^2$ : $T \times_S T \simeq G \times T$. La même correction vaut plus
bas pour $T \times_S T \times_S T \simeq G \times G \times T$, que la page écrit
« $\Delta_3 T \times (G \times G \times G)$ », alors qu'elle indexe ensuite
correctement les composantes par $(g, g')$.} On a
\[
  \varphi_g^{-1}\bigl(U_g \cap (T' \times_S T')\bigr) = g^{-1}(T') \cap T' ,
\]
qui vaut $T'$ si $g = e$ et $T''$ si $g \neq e$. Ainsi $T' \times_S T'$ est
somme de sa diagonale, isomorphe à $T'$, et d'ouverts $U'_g$ ($g \neq e$)
isomorphes à $T''$.

\emph{Le produit triple.} De même $T \times_S T \times_S T = \coprod_{(g, g')}
V_{g,g'}$ avec $V_{g,g'} = (g \times g' \times e)(\Delta_3 T) \simeq T$, et
$V'_{g,g'} = V_{g,g'} \cap T'^{3}$ s'identifie à $g^{-1}(T') \cap g'^{-1}(T')
\cap T'$, qui est $T''$ sauf pour $g = g' = e$, où c'est la diagonale $\simeq
T'$.

\emph{Données de descente.} Une donnée de descente sur un objet $X$ au-dessus
de $T'$, relativement à $T' \to S$, est un isomorphisme $p_1^{*}X \simeq
p_2^{*}X$ sur $T' \times_S T'$, c'est-à-dire un système d'isomorphismes sur les
composantes : sur la diagonale, c'est l'identité, et pour $g \neq e$, c'est un
isomorphisme $X|_{T''} \simeq g^{*}(X|_{T''})$. La condition de cocycle est
automatique sur la diagonale $V'_{e,e}$ ; sur les autres composantes, elle ne
porte que sur $T''$, et c'est exactement la condition de cocycle d'une donnée
de descente sur $X|_{T''}$ relativement au $G$-torseur $T'' \to S''$. D'où, comme
l'écrit la page en le soulignant :

\begin{quote}
\emph{une donnée de descente sur $X$ relativement à $T' \to S$ est la même
chose qu'une donnée de descente sur $X|_{T''}$ relativement à $T'' \to S''$,
c'est-à-dire qu'une action de $G$ sur $X|_{T''}$ compatible à son action sur
$T''$.}
\end{quote}

Si par exemple $X$ est projectif sur $T'$, la descente galoisienne le long du
revêtement fini $T'' \to S''$ est effective, et la donnée revient à un
isomorphisme $X|_{T''} \simeq X_0 \times_{S''} T''$, avec $X_0$ projectif sur
$S''$.\footnote{La page dit « proj. » des deux côtés. Ce que garantit la
descente galoisienne des schémas quasi-projectifs, c'est un $X_0$ propre sur
$S''$ muni d'un faisceau ample, déduit par norme d'un faisceau ample de
$X|_{T''}$ ; c'est le sens qu'on donne ici à « projectif ».}

\emph{La question d'effectivité} (page 34) devient donc la suivante : étant
donné $X_0$ sur $S - \{s\}$, et un prolongement $X$ de $X_0 \times_{S - \{s\}}
(T - f^{-1}(s))$ en le seul point $t$ de $f^{-1}(s)$ — donc au-dessus de $T' =
(T - f^{-1}(s)) \cup \{t\}$ —, par exemple la restriction d'un prolongement
sur $T$ tout entier, montrer que $X$ est isomorphe à l'image réciproque d'un
schéma sur $S$ prolongeant $X_0$.\footnote{La page écrit d'abord $X'_0$, puis
$X_0$ ; il s'agit du même objet.} Le dossier s'arrête sur cette question, sans
y répondre.

On sait que la descente des schémas le long d'un morphisme fpqc n'est pas
effective en général, et qu'elle l'est pour les schémas quasi-projectifs
munis d'un faisceau ample compatible à la donnée de descente (SGA~1, exposé
VIII, § 7, cité de mémoire) ; l'exemple de la page est construit pour isoler
ce qui manque quand le morphisme n'est pas fini : tout se passe au-dessus de
$S''$, où le revêtement est fini, sauf en un point, où il n'y a plus de
donnée de descente du tout, et où la question devient une question de
prolongement.

\end{document}
